\section[Keterkaitan Buku Ajar dengan RPS OBE]{Keterkaitan Buku Ajar dengan RPS\\Berbasis OBE}

\subsection{Alignment dengan Sub-CPMK}

Outcome-Based Education (OBE) menekankan organisasi seluruh sistem pembelajaran di sekitar kompetensi yang harus dicapai mahasiswa pada akhir pengalaman belajar \cite{spady1994,obe_springer2024}. Menurut Spady, fokus bergeser dari input (materi yang diajarkan) ke output (kemampuan yang terdemokratisasi), sehingga asesmen dan strategi pembelajaran harus selaras dengan outcome yang telah didefinisikan. Kerangka OBE yang dirancang Noushad dan sumber lain menekankan pentingnya alignment antara learning outcomes, materi ajar, metode pembelajaran, dan instrumen evaluasi.

Rencana Pembelajaran Semester (RPS) berbasis OBE merupakan dokumen yang memetakan capaian pembelajaran ke aktivitas dan asesmen secara eksplisit. Buku ajar ini berfungsi sebagai sumber belajar utama yang menyediakan materi, contoh, latihan, dan instrumen evaluasi yang selaras dengan RPS tersebut. Setiap bab dalam buku ini dikaitkan langsung dengan Sub-CPMK tertentu, memungkinkan dosen dan mahasiswa menilai pencapaian kompetensi secara terukur.

Pemetaan bab terhadap Sub-CPMK mata kuliah Logika Matematika adalah sebagai berikut:

\begin{itemize}
  \item \textbf{Bab III} mencakup Sub-CPMK 1.1, 1.2, 1.3 (Proposisi dan Operator Logika)
  \item \textbf{Bab IV} mencakup Sub-CPMK 2.1, 2.2, 2.3 (Ekuivalensi dan Tautologi)
  \item \textbf{Bab V} mencakup Sub-CPMK 3.1 (Logika Predikat)
  \item \textbf{Bab VI} mencakup Sub-CPMK 3.2, 4.1 (Aturan Inferensi dan Metode Bukti)
  \item \textbf{Bab VII} mencakup Sub-CPMK 4.2 (Induksi Matematika)
  \item \textbf{Bab VIII} mencakup Sub-CPMK 4.3 (Prolog)
\end{itemize}

\subsection{Peran Buku Ajar terhadap RPS}

Materi, aktivitas pembelajaran, latihan, dan asesmen dalam setiap bab dirancang untuk mengukur pencapaian Sub-CPMK yang bersangkutan \cite{in4obe}. Dengan demikian, buku ajar ini tidak sekadar menyajikan konten, melainkan berfungsi sebagai alat bantu pencapaian kompetensi yang terintegrasi dengan RPS dan kurikulum Program Studi Teknik Informatika.
